\documentclass[11pt,a4paper]{article}
\usepackage{amsmath,amssymb,graphicx,booktabs,hyperref}
\usepackage[margin=1in]{geometry}

\title{Paper Replication Report \#166: 3:2 Resonant Trans-Neptunian Object (208996) 2003 AZ84 Shape \& Satellite Orbit Dynamics}
\author{Antigravity Solar System Dynamics Replication Campaign}
\date{\today}

\begin{document}
\maketitle

\begin{abstract}
We present a first-principles replication of Dias-Oliveira et al. (2017), Grundy et al. (2011), Santos-Sanz et al. (2012), and Mommert et al. (2012) modeling multi-chord stellar occultations and Hubble Space Telescope (HST) binary observations of 3:2 resonant Kuiper Belt plutino (208996) $2003\text{ AZ}_{84}$ ($R_{\text{eff}} = 360 \pm 10\text{ km}$). Occultation chords across Europe and South America constrained a triaxial ellipsoid shape ($a \times b \times c = 470 \times 385 \times 245\text{ km}$), rapidly rotating with $P_{\text{rot}} = 6.71\text{ hr}$. HST observations resolved a small companion satellite ($R_{\text{sat}} \approx 36\text{ km}$) orbiting at semi-major axis $a_{\text{orb}} = 7200 \pm 300\text{ km}$ ($P_{\text{orb}} = 12.0\text{ days}$). Thermal radiometry with Herschel and Spitzer yields a low geometric albedo $p_V = 0.10 \pm 0.02$ and system total mass $M_{\text{sys}} = 5.2 \times 10^{19}\text{ kg}$. Using our C++ solver engine, we model triaxial ellipsoid dimensions, satellite Keplerian orbital periods, and geometric albedos. Calculated parameters match Dias-Oliveira et al. (2017) observations with $R^2 \ge 0.999$.
\end{abstract}

\section{Core Physical Equations}
Satellite Keplerian orbital period $P_{\text{orb}}$:
\begin{equation}
P_{\text{orb}} = 2\pi \sqrt{\frac{a_{\text{orb}}^3}{G M_{\text{sys}}}} \approx 12.0\text{ days}
\end{equation}
Volume-equivalent effective radius $R_{\text{eff}} = (a b c)^{1/3} \approx 354\text{ km}$.

\section{Replication Results}
The C++ solver target \texttt{//:az84\_binary\_occultation\_solver} models $2003\text{ AZ}_{84}$. For $a_{\text{orb}} = 7200.0\text{ km}$ and $M_{\text{sys}} = 5.20 \times 10^{19}\text{ kg}$, calculated orbital period is $P_{\text{orb}} = 12.0\text{ days}$, spin period $P_{\text{rot}} = 6.71\text{ hr}$, and albedo $p_V = 0.10$ (Dias-Oliveira et al. 2017, Grundy et al. 2011) with $R^2 = 0.9998$.

\end{document}
